\section{Условия линейной сходимости}\label{sec:linear}

В общей монотонной задаче оценка \eqref{eq:hpe-energy} даёт скорость
зазора $\mathcal O(1/S)$ по накопленному весу. Чтобы получить
геометрическое убывание, нужны условия, при которых вычисленная
невязка служит индикатором расстояния до решения.

\subsection{Оценка расстояния через невязку}

Обозначим $T=F+N_C$, $Z=T^{-1}(0)$ и
$d(x)=\dist(x,Z)=\inf_{z\in Z}\norm{x-z}$.
Предположим, что на принятых пробах нижнего порога выполняется
\begin{equation}\label{eq:error-bound}
 d(u)\le\kappa\,\dist(0,T(u)),\qquad \kappa>0.
\end{equation}
Название этого локального свойства --- метрическая
субрегулярность~\cite{leventhal}. Оно позволяет судить по малой
невязке о близости к множеству решений.
Из \eqref{eq:error-bound} следует $d(u)\le\kappa\norm v$
для любого вычисленного $v\in T(u)$.

Если $F$ $\mu$-сильно монотонен на $C$, то
$\ip{v}{u-z^\star}\ge\mu\norm{u-z^\star}^2$ по монотонности
нормального конуса, и можно взять $\kappa=1/\mu$.
Оценка \eqref{eq:error-bound} также покрывает билинейные игры
без сильной монотонности. Например, для $C=\R^{2d}$,
$F(x,y)=(By,-B^\top x)$ и квадратной невырожденной $B$ имеем
$Z=\{0\}$ и $\norm{F(x,y)}\ge\sigma_{\min}(B)\norm{(x,y)}$.
При этом $\ip{F(x,y)}{(x,y)}=0$ для всех $x,y$.
Именно связь остатка с расстоянием здесь обеспечивает сжатие.

\subsection{Сжатие и явный критерий остановки}

Меняем нижний порог на $\Gamma_0=1/\kappa$ и на нём
останавливаемся при $\kappa\norm v\le\eta$.
В алгоритме~\ref{alg:hpe} этот критерий заменяет остановку
по сумме весов. На выход возвращаем принятую точку $u$.
Пусть $\sigma=1/21$.
Из \eqref{eq:hpe-energy}, \eqref{eq:hpe-res} и монотонности
получаем для всех принятых шагов
\begin{equation}\label{eq:linear-energy}
 d(x^+)^2\le d(x)^2-(1-\sigma^2)\norm{u-x}^2,\qquad
 \norm v\le\Gamma(1+\sigma)\norm{u-x}.
\end{equation}
Для первого неравенства подставляем ближайшее к $x$ решение.
Для второго используем треугольник в определении $e_H$.
На нижнем пороге \eqref{eq:error-bound} даёт
$d(x)\le d(u)+\norm{u-x}\le(2+\sigma)\norm{u-x}$.
Следовательно,
\begin{equation}\label{eq:hpe-q}
 d(x^+)^2\le q\,d(x)^2,\qquad
 q=1-\frac{1-\sigma^2}{(2+\sigma)^2}
   =\frac{1409}{1849}.
\end{equation}
Это коэффициент для \emph{квадрата} расстояния.
Каждый шаг с $\Gamma>\Gamma_0$ проходит проверку по длине и потому
уменьшает квадрат расстояния как минимум на
$c=(440/441)(29r/128)^2$. Таких шагов за весь запуск не больше
$B=\lfloor R^2/c\rfloor$.

На пороге те же неравенства дают
$(\kappa\norm v)^2\le(11/10)d(x)^2$.
Положим
\[
  m=\ceil{\log_{1/q}\max\{1,(11/10)R^2/\eta^2\}}.
\]
В префиксе из $B+m+1$ шагов есть проба на нижнем пороге,
до которой уже произошло не менее $m$ сжатий.
На этой пробе $\kappa\norm v\le\eta$ и $d(u)\le\eta$.
Поэтому число внешних шагов имеет порядок
\begin{equation}\label{eq:linear-steps}
 \mathcal O\!\left(1+R^2/r^2+\log(1+R/\eta)\right).
\end{equation}
Если \eqref{eq:error-bound} выполнено только в
$\ball_a(z^\star)\cap C$, достаточно старта в $\ball_{a/3}(z^\star)$.
Расстояние внешних центров до этого решения не растёт, а
\eqref{eq:linear-energy} в виде оценки относительно конкретного
$z^\star$ даёт $\norm{u-x}\le2\norm{x-z^\star}$.
Значит, все следующие принятые точки лежат в $\ball_a(z^\star)$.

\subsection{Бюджеты и случай сильной монотонности}

Положим $\mathcal K_\eta=1+R^2/r^2+\log(1+R/\eta)$.
Во всех принятых шагах $\Gamma\ge1/\kappa$.
Из локальной оценки \eqref{eq:hpe-curv} и подсчёта проб в
разделе~\ref{sec:hpe-cost} получаем
\begin{align}
 N_{\nabla p_i}
 &=\OO\!\left(17^{2n-1}\mathcal K_\eta
               [1+\sqrt{\kappa A_i}+\sqrt{L_{1,i}R}]\right),
               \label{eq:linear-grad}\\
 N_{Q_i}
 &=\OO\!\left(17^{2n-1}\mathcal K_\eta[1+\kappa M_i]\right).
               \label{eq:linear-op}
\end{align}
Число вызовов зависит от точности логарифмически при фиксированных
параметрах задачи. Масштаб оператора остаётся линейным по $\kappa M_i$,
а градиентный масштаб растёт как $\sqrt{\kappa A_i}$.

При $\mu$-сильной монотонности на пороге $\Gamma_0=\mu$ можно
оставить в \eqref{eq:hpe-energy} член
$2\Gamma_0^{-1}\ip{v}{u-z^\star}\ge2\norm{u-z^\star}^2$.
Сумма $2\norm{u-z^\star}^2+(440/441)\norm{u-x}^2$ не меньше
$(440/661)\norm{x-z^\star}^2$ по взвешенному неравенству
Коши--Буняковского. Поэтому коэффициент в \eqref{eq:hpe-q}
улучшается до $221/661$.

В этом случае вычисленная невязка также сертифицирует
абсолютный зазор на всём $C$. По выпуклости, монотонности
и интегрированию $F$ вдоль отрезка $[w,u]$ имеем
\[
 G(u,w)\le \ip{v}{u-w}-\frac\mu2\norm{u-w}^2
          \le\frac{\norm v^2}{2\mu}\qquad(w\in C).
\]
В первом неравенстве $Q(w)$ заменяется интегралом верхней оценки
по отрезку, а сильная монотонность даёт вычитаемый квадратичный
член. Следовательно, достаточно потребовать
$\norm v\le\sqrt{2\mu\varepsilon}$ для абсолютного зазора
$\varepsilon$.

\subsection{Перезапуски пути А при сильной выпуклости}

Есть и прямое следствие для пути А. Если сумма $p$ является
$\mu$-сильно выпуклой на $C$, то по условию ВН
$G(y,z^\star)\ge(\mu/2)\norm{y-z^\star}^2$.
Теорема~\ref{thm:patha} при $\rho=\mu/8$ даёт
\[
 \frac\mu2\norm{z_{\rm out}-z^\star}^2
 \le G(z_{\rm out},z^\star)
 \le\frac\mu8\norm{x_0-z^\star}^2 .
\]
Один запуск уменьшает расстояние до решения вдвое.
Повторение из последнего выхода даёт коэффициент $2^{-j}$
после $j$ запусков. Достаточно сильной выпуклости суммы функций.
Цена каждого запуска по-прежнему определяется общим масштабом
градиента и оператора внутри пары из \eqref{eq:pa-cost}.
